\begin{enumerate}[label=\textbf{\ansnum.\arabic*}]
  \item $202=128+64+8+2=\texttt{0b11001010}=\texttt{0xca}$ ($12\cdot16+10$). $\texttt{0xff}=255$, $\texttt{0b101101}=32+8+4+1=45$.
  \item $2^n$ values. $2^{10}=1024\ge1000$, so $10$ bits.
  \item $1280\cdot720=921\,600$ bytes per image, $\cdot60=55.3$ MB/s $=442$ Mbit/s. No: it is more than four times the $100$ Mbit/s of the link (even gigabit Ethernet would be half used); compress the video or reduce size and rate.
  \item $60\cdot10^6\cdot4=240$ MB, $120$ MB, $60$ MB. None fits into $32$ MB; it would need a network of fewer than $32\cdot10^6$ parameters in \texttt{int8}.
  \item $0.1$ and $0.2$ have no exact binary representation, the stored values are slightly off and the sum is $0.30000000000000004$; use \texttt{math.isclose}. A \texttt{float32} holds about $7$ decimal digits ($23$ mantissa bits).
  \item Registers, L1, L2, L3 cache, RAM, SSD (and network, and disk). The faster a memory is, the more expensive per byte and the smaller (and the more power it needs), so a hierarchy gives most of the speed at a fair price: frequently used data stays in the small fast levels.
  \item A Python loop executes interpreter instructions for each element and the Python objects lie scattered in memory (poor use of cache lines); NumPy runs one compiled loop over a contiguous block of memory, which streams through the cache and can use vector instructions of the CPU.
  \item RAM is fast, volatile working memory (content lost without power); an SSD is slower, larger, non-volatile storage (content stays). A program must be loaded from the SSD into RAM to run.
  \item A GPU has thousands of simple cores that perform the same operation on many data at once (matrix multiplication); a CPU has few powerful cores with good branch prediction, optimised for decisions and varied sequential work.
  \item[\textbf{\ansnum.10}] (1) Power: LEDs on the board, battery voltage, a different charger or cable. (2) Firmware: a screen or beep codes? Connect a monitor and keyboard; a serial console (UART) shows firmware output. (3) Bootloader/disk: does a boot menu appear? Boot from a USB stick with a live system and look at the disk with \texttt{lsblk} and \texttt{fsck}. (4) Kernel: read the messages on the screen or, after a successful boot of the live system, \verb!journalctl -b -1! of the failed boot; \texttt{dmesg} shows driver errors. (5) Services/network: \verb!systemctl --failed!, \verb!ip a!, \verb!journalctl -u NAME!. Move one step at a time and note which stage is the last that works.
\end{enumerate}
